\chapter{Introducción}
\label{anexo1:cap:introduccion}
En este documento se recoge la especificación de los requisitos del sistema, que incluye los
requisitos funcionales, no funcionales y de información, así como los casos de uso, distintos
actores involucrados y, finalmente, la matriz de trazabilidad de requisitos. La estructura de
este documento sigue la descrita por \fullcite{DuranBernardez}.

Estos requisitos se han obtenido a partir de los apartados de \enquote{descripción} y
\enquote{objetivos funcionales} de la propuesta del trabajo de fin de grado tal y como se recoge
en la página web de la asignatura\footnote{\url{https://diaweb.usal.es/diaweb/index.jsp}}. La
descripción y objetivos funcionales dicen lo siguiente:
\begin{quote}
      \textbf{Descripción:}

      El proyecto consiste en desarrollar una plataforma de entretenimiento interactivo basada en la
      estimación de posturas corporales, utilizando tecnología de visión por computadora. La
      plataforma permitirá a los usuarios participar en una variedad de juegos y actividades que se
      controlan mediante los movimientos del cuerpo, sin necesidad de controles físicos adicionales.
      A través del uso de cámaras (webcams o cámaras de dispositivos móviles), el sistema detectará
      y analizará la postura y los gestos del usuario en tiempo real, permitiendo una experiencia de
      juego inmersiva y dinámica.

      \textbf{Objetivos funcionales:}

      \begin{itemize}
            \item Implementar un sistema de detección y seguimiento de posturas en tiempo real a través
                  de cámaras.
            \item Desarrollar una variedad de juegos y actividades interactivas que utilicen la postura
                  del cuerpo como control principal.
            \item Ofrecer personalización de actividades según las capacidades físicas del usuario,
                  adaptando el nivel de dificultad.
            \item Integrar una interfaz intuitiva y amigable que permita a los usuarios seleccionar y
                  personalizar las experiencias de entretenimiento.
            \item Proporcionar métricas y estadísticas de rendimiento (movimiento, esfuerzo, etc.) que
                  promuevan el ejercicio físico y la mejora personal.
      \end{itemize}
\end{quote}

